% !TEX program = xelatex
\documentclass[10pt, a4paper]{article}

% ===== Page setup =====
\usepackage[left=1.0cm, right=1.0cm, top=0.6cm, bottom=0.6cm]{geometry}
\usepackage{titlesec}
\usepackage{enumitem}
\usepackage{tabularx}
\usepackage{xcolor}
\usepackage{hyperref}
\usepackage{fancyhdr}
\usepackage{fontspec}
\usepackage{xeCJK}

% ===== Fonts =====
\setCJKmainfont{Songti SC}
\setCJKsansfont{Hiragino Sans GB}
\setCJKmonofont{Hiragino Sans GB}

% ===== Colors =====
\definecolor{primary}{RGB}{0, 90, 160}
\definecolor{darktext}{RGB}{40, 40, 40}
\definecolor{graytext}{RGB}{100, 100, 100}

% ===== Hyperlinks =====
\hypersetup{
    colorlinks=true,
    linkcolor=primary,
    urlcolor=primary,
    pdfborder={0 0 0}
}

% ===== No page numbers =====
\pagestyle{empty}

% ===== Paragraph settings =====
\setlength{\parindent}{0pt}
\setlength{\parskip}{0pt}

% ===== Section title style =====
\titleformat{\section}
    {\large\bfseries\color{primary}}
    {}
    {0em}
    {}
    [\vspace{-0.6em}\textcolor{primary}{\rule{\textwidth}{0.8pt}}]
\titlespacing*{\section}{0pt}{0.35em}{0.1em}

% ===== List style =====
\setlist[itemize]{
    leftmargin=1.5em,
    itemsep=0pt,
    parsep=0pt,
    topsep=0em,
    label=\textcolor{primary}{\textbullet}
}

% ===== Custom commands =====
% Entry: title | right-side information
\newcommand{\entry}[2]{%
    \vspace{0.1em}
    \noindent\textbf{#1}\hfill{\small\color{graytext}#2}\\[-0.4em]
}

% Subtitle row
\newcommand{\entrysub}[2]{%
    \noindent{\small #1}\hfill{\small\color{graytext}#2}\\[-0.5em]
}

\begin{document}

% ===== Header: name and contact information =====
\begin{center}
    {\LARGE\bfseries 袁毅}\quad {\Large Yi Yuan}\\[0.3em]
    {\small
        \textbf{邮箱：} yuanyithu@126.com \quad
        \textbf{GitHub：} \href{https://github.com/yuanyithu}{github.com/yuanyithu} \\[0.2em]
        男\enspace|\enspace 清华大学物理系博士研究生
    }
\end{center}

% ===== Education =====
\section{教育背景}

\entry{清华大学 \enspace 物理学博士}{2023.08 -- 2028.06（预计）}
\entrysub{物理系 \enspace 博士资格考试现代物理方向排名\textbf{第 1} \enspace 量子计算机原理与架构（A+）}{北京}

\entry{清华大学 \enspace 数理基础科学学士}{2019.08 -- 2023.06}
\entrysub{物理系 \enspace GPA：3.89/4.0（排名：18/116）\enspace 优秀毕业生 \enspace 叶企孙奖提名}{北京}

% ===== Skills =====
\section{技术能力}
\begin{itemize}
    \item \textbf{编程语言：} Python（日常科研与开发的主要语言）
    \item \textbf{AI/机器学习：} 熟悉 Transformer 架构与 RLHF 流程（PPO/GRPO）；熟练使用 JAX 和 PyTorch；具有 CNN 建模与 LLM API 应用开发经验
    \item \textbf{数据与计算：} NumPy 向量化计算、多进程、数据驱动建模与回测系统开发
    \item \textbf{工具：} Git、Linux 服务器 \quad \textbf{语言：} 英语，可熟练阅读与写作科学文献
\end{itemize}

% ===== Publications =====
\section{论文发表}
\noindent{\small \textbf{Yi Yuan}, Yuanchen Zhao, and Dong E. Liu. ``Zeno-Enhanced Probabilistic Error Cancellation with Quantum Error Detection Codes.'' \textit{arXiv preprint arXiv:2605.12149}, 2026. \href{https://arxiv.org/abs/2605.12149}{arXiv}}\\[-0.2em]
\noindent{\small 提出面向近期含噪量子设备的微扰型 QEDC+PEC 协议，通过后选择重塑有效噪声信道，并利用 PEC 消除剩余逻辑错误。}

% ===== Research and projects =====
\section{科研与项目}

\entry{量子计算与量子信息研究}{2022.09 -- 至今}
\entrysub{清华大学物理系}{Python, JAX, Tensor Network}
\begin{itemize}
    \item \textbf{Zeno 增强的 QEDC+PEC 协议}（2024 -- 2026，arXiv 预印本）：结合量子错误探测码与概率误差抵消，面向近期含噪量子设备开展误差缓解；在 200 比特 Iceberg-code 逻辑 GHZ 制备中，相比标准 PEC 将采样开销降低 \textbf{3--4 个数量级}，同时保持 $F\simeq0.956$
    \item \textbf{经典影子误差缓解}（2024）：使用张量网络方法计算含噪量子线路中经典影子的误差缓解复杂度，并分析其随系统规模的标度行为
    \item \textbf{本科毕业论文：量子退火复杂度}（2022.09 -- 2023.06）：基于虚时间时空瞬子作用量，解释双阱势基态能量搜索中量子退火与经典 QMC 复杂度相同的原因，并指出多路径量子加速的可能性
\end{itemize}

\entry{代表项目}{}
\begin{itemize}
    \item \textbf{SEVI 粒子散射分析：} 基于 PyTorch 构建 CNN 拟合散射振幅参数，在不注入先验知识的情况下达到解析方法精度
    \item \textbf{LLM 论文总结工具}（2025）：自动获取每日 arXiv 论文，通过 Kimi API 生成总结，并完成本地部署
\end{itemize}

% ===== Honors and leadership =====
\section{荣誉与工作经历}
\noindent{\small 荣誉：2021 年清华大学校友 -- 张明为奖学金 \quad 2020--2021 年多次获得物理系学业优秀、创新创业与综合优秀奖学金}\\[-0.2em]
\noindent{\small 工作经历：2023.08--2025.01 清华大学物理系学生辅导员 \quad 2019--2020 团支部书记（所在支部获评“甲级团支部”）\quad 2024--2025 优秀学生干部}

\end{document}
